% ---- Modele TP PC V2 -----------
\documentclass[a4paper,9pt,french]{scrartcl}
\usepackage[T1]{fontenc}%---pour l'encodage d'entr\'ee
\usepackage{array}%---pour les tableaux
\usepackage{titlesec}%---modification des titres
\usepackage{ulem}%---pour le soulignage
\usepackage[dvipsnames]{xcolor}%---pour avoir d'autres couleurs
\usepackage{color}%---pour avoir les couleurs de base
\usepackage{babel}%---pour les trucs de langue et la francisation

\usepackage{graphicx}%----pour mettre des images
\usepackage[utf8]{inputenc}%---encodage
\usepackage{geometry}%---pour modifier les tailles et mettre a4paper
%\usepackage{awesomebox}%---pour les boites d'exercices, de pbq et de croquis ---d\'esactiv\'e pour les TP de PC
\usepackage{tikz}%---pour dessiner + d\'ependance de chemfig
%\usepackage{tabularx}%---pour dimensionner automatiquement les tableaux avec variable X
\usepackage{fancyhdr}%---pour les en-t\^ete personnalis\'ees
\usepackage{blindtext}%---pour les liens
\usepackage{hyperref}%---pour les liens (\`a mettre en dernier)
\usepackage{caption}%---pour la francisation de la l\'egende table vers Tableau
 \usepackage{float}% --- pour placer les figures et tableaux où l'on souhaite avec argument H
 \usepackage{multicol}
 \usepackage{amsmath}
 \usepackage{geometry}
 \geometry{vmargin=2.5cm}

\captionsetup{labelfont=sc}%---pour la francisation de la l\'egende table vers Tableau
\def\frenchtablename{Tableau}%---pour la francisation de la l\'egende table vers Tableau
\pagestyle{fancy}
\renewcommand\headrulewidth{1pt}
\fancyhead[L]{TP 1}
\fancyhead[R]{Saux Paulhenry}
\title{Compte rendu de TP 6}
\date{}
\author{}

\geometry{a4paper}
% ----------- Création des commandes de couleur ----------
\makeatletter
\newcommand{\coloruline}[2]{%
        \UL@protected\def\temp@uline{\leavevmode \bgroup
            \markoverwith{\textcolor{#1}{\rule[-0.5ex]{2pt}{0.4pt}}}%
        \ULon}%
        \temp@uline{#2}%
    }

    \newcommand{\coloruuline}[2]{%
        \UL@protected\def\temp@uuline{\leavevmode \bgroup
            \UL@setULdepth
            \ifx\UL@on\UL@onin \advance\ULdepth2.8\p@\fi
            \markoverwith{\textcolor{#1}{\lower\ULdepth\hbox
                {\kern-.03em\vbox{\hrule width.2em\kern1\p@\hrule}\kern-.03em}}}%
        \ULon}%
        \temp@uuline{#2}%
    }

\makeatother

\renewcommand{\thesection}{\arabic{section}}
\renewcommand{\thesubsection}{\arabic{subsection}}
\renewcommand{\thesubsubsection}{\arabic{subsubsection}}
\begin{document}
% ---------------- Modification des titres de niveau 1,2 et 3 --------------------
\titleformat
{\section} % command
%[display] % shape
{\Large} % format
{\coloruuline{red}{\thesection. }} % label
{0ex} % sep
{\coloruuline{red}} % before-code
[] % after-code

\titleformat
{\paragraph} % command
%[display] % shape
{\Large} % format
{\coloruline{NavyBlue}{\theparagraph. }} % label
{0ex} % sep
{\coloruline{NavyBlue}} % before-code
[] % after-code

\titleformat
{\subsection} % command
%[display] % shape
{\Large} % format
{\coloruline{red}{\thesection.\thesubsection. }} % label
{0ex} % sep
{\coloruline{red}} % before-code
[] % after-code

\titleformat
{\subsubsection} % command
%[display] % shape
{\large} % format
{\coloruline{ForestGreen}{\thesection.\thesubsection.\thesubsubsection. }} % label
{0ex} % sep
{\coloruline{ForestGreen}} % before-code
[] % after-code

\titleformat
{\chapter} % command
%[display] % shape
{\large} % format
{\coloruline{NavyBlue}{\thechapter \:- }} % label
{0ex} % sep
{\coloruline{NavyBlue}} % before-code
[] % after-code


% ---------------- Début du corps du document ------------------------------------

\section{Introduction}
\subsection{Objectifs}
\subsection{Sommaire}
Ce TP se décompose en quatre parties :
\begin{enumerate}
 \item
\end{enumerate}
\section{Protocole}
\subsection{Matériel}
\begin{itemize}
\end{itemize}
\subsection{Description du montage expérimental}
\begin{center}
%\includegraphics[scale=0.3]{circuit1}
\end{center}
\subsection{Mesures à effectuer}
\section{Mesures}
\subsection{Données}

\subsection{Incertitudes}

\section{Graphiques}
\subsection{Analyse de la tension aux bornes de la résistance \& Étude de la résonnance en intensité}
% \begin{figure}[H]
%  \begin{center}
% \includegraphics[scale=0.4]{1}
%  \end{center}
% \caption{Diagramme de Bode en amplitude}
% \end{figure}

\section{Exploitation des résultats}
\section{Conclusion}

\end{document}
